\chapter{Background and Problem Setting}
\label{sec:background}

\section{Reverse Engineering as an Evidence-Gathering Problem}

Reverse engineering is not a single question answered by a single tool. Analysts move among disassembly, decompilation, string inspection, structural triage, pattern matching, and selective execution in order to construct an evidence-backed explanation of program behavior. This makes the domain a strong candidate for studying tool-orchestrated agents.

\begin{itemize}
\item Reverse engineering tasks often require both broad survey steps and narrow recovery steps.
\item Malware-oriented questions raise the additional need for analyst caution, provenance tracking, and explicit uncertainty handling.
\item The quality of an answer depends not only on the final narrative but also on whether the system gathered the right supporting artifacts.
\end{itemize}

\begin{itemize}
\item \TODO{Define the task categories used in this thesis, such as default analysis, configuration recovery, behavior recovery, packing triage, and anti-analysis investigation.}
\item \TODO{State what counts as acceptable evidence for binary understanding claims in this thesis.}
\end{itemize}

\section{LLM Agents, Tool Use, and Orchestration}

The system studied here is explicitly modular. Rather than expecting one model invocation to reason, search, inspect artifacts, and self-critique all at once, the workflow decomposes the task into stages with different responsibilities. This chapter should explain why that decomposition is appropriate for reverse engineering.

\begin{itemize}
\item Planner stages decompose the query and allocate work.
\item Worker stages gather concrete evidence through tool use.
\item Validator stages review coverage, evidence quality, and missing artifacts.
\item Reporter stages synthesize the final answer for the run artifact.
\end{itemize}

\begin{itemize}
\item \TODO{Explain the intuition behind planner-worker-validator-reporter style decomposition.}
\item \TODO{Clarify how this differs from chain-of-thought prompting or single-agent tool use.}
\item \TODO{Discuss why explicit stage contracts matter when tools can produce noisy or partial evidence.}
\end{itemize}

\section{MCP-Style Tool Integration}

The repository uses MCP-style servers to expose tool capabilities to the runtime. This matters because the thesis argument is partly about interface design: how tool affordances are surfaced to agents, how outputs are normalized, and how evidence is routed back into shared state.

\begin{itemize}
\item Core static lane: Ghidra, strings, FLOSS, capability analysis, hash lookups, YARA, binwalk, and related enrichers.
\item Optional dynamic or execution-oriented lane: VM and sandbox-style tooling, packet or process observation, and task-specific execution probes.
\item Artifact management layer: bundles, manifests, run records, judge outputs, and experiment aggregates.
\end{itemize}

\begin{itemize}
\item \TODO{Document how MCP servers are loaded and partitioned into static versus dynamic tool groups.}
\item \TODO{Explain the tradeoff between uniform tool access and tool-specific prompting or guardrails.}
\item \TODO{Add a diagram showing how tool requests and results move through shared state.}
\end{itemize}

\section{Binary Analysis Tooling in Scope}

This thesis should also introduce the technical role of each major tool class so later architecture chapters can assume that context.

\begin{itemize}
\item Ghidra provides decompilation, disassembly, symbol context, and structural inspection.
\item String and FLOSS-style tooling surface literal and decoded text evidence.
\item Capability and signature tools provide fast high-level hints that can be corroborated later.
\item Dynamic or VM-backed tooling is reserved for cases where static evidence is insufficient or ambiguous.
\end{itemize}

\begin{itemize}
\item \TODO{Explain why Ghidra is the primary structural analysis backend in this system.}
\item \TODO{Clarify which tasks can be completed from static evidence alone and which ones benefit from dynamic execution.}
\item \TODO{Describe how UPX, packing, and unpacking are treated as preprocessing versus analysis-time decisions.}
\end{itemize}

\section{Problem Formulation for the Thesis}

At a high level, each experiment asks whether a workflow configuration can answer a binary-analysis task better than a baseline. In the maintained harness, the atomic evaluation unit is a sample-task pair. That framing should become the formal problem statement for the evaluation chapter.

\begin{itemize}
\item Inputs: binary sample, task prompt, tool availability, workflow configuration, and model configuration.
\item Outputs: final analysis report, intermediate evidence artifacts, judge score, and run metadata.
\item Comparison axis: score deltas, task success rates, error rates, and qualitative failure analysis across matched tasks.
\end{itemize}

\begin{itemize}
\item \TODO{Convert this into a concise formal problem statement that the rest of the thesis can reference.}
\item \TODO{State what a ``weaker baseline'' means in the context of this system and why that baseline is fair.}
\end{itemize}
